\section{Benchmark design}
\label{sec:design}

\subsection{Two tracks}

\paragraph{Service track (S).}
Tasks depend on mock services exposed over MCP: ticketing, an inventory database, a deploy API and documentation search.
A fault proxy sits between the agent and the services on the task's internal network and runs in its own container, so the agent sees it only through its MCP interface.
Every harness receives the same bytes for the same faulted call.
This is the primary cross-harness board.

\paragraph{Native track (N).}
Repository tasks rely on the harness's built-in tools: read, write, edit, search and shell.
An injector loaded through each harness's own extension point applies the faults: a Claude Code plugin, Codex command hooks, a Pi extension, an OpenCode plugin, and the hooks of fxh.
Harnesses differ in what their extension points can change, so native-track results are compared only within a parity set: the fault types that a conformance suite shows each harness in the comparison delivers to the model (Section~\ref{sec:conformance}).
fx~0.0.13 has no interception point for its built-in tools, so its native-track row is fxh, a fork of fx~0.0.13 maintained by the author that adds tool hooks and changes nothing else in the agent loop.
The fxh row measures that fork; it is not a measurement of stock fx on the native track.

\subsection{What the fixed-model condition controls}
\label{sec:fixed}

Every configuration runs the same main model, \code{claude-sonnet-5-5}, with the same nominal reasoning setting (\code{medium}), in full autonomy inside the sandbox (Table~\ref{tab:configs}).
Tasks, profiles, seed, attempt indices and container images are also shared.
The condition does not equalize everything the harness sends to the model.
Recorded differences are: Pi~1.0.3 sends a top-level reasoning effort \code{high} with a mid-conversation effort \code{medium}; Pi and OpenCode request summarized thinking while Claude Code requests it omitted; OpenCode sends one request per run to a smaller model (\code{claude-haiku-5-5}) to title the session; fx and fxh reach the model through the Vercel AI Gateway instead of the provider API; and Codex CLI, which speaks only the Responses API, reaches the model through a bridge we wrote that translates its requests to the Messages API.
We treat these as part of each harness configuration and interpret differences between rows at the configuration level (harness, version, tool interface and model route), not as the causal effect of one isolated harness mechanism.
The specification also defines a native axis, where each harness runs its vendor's model; it was not run in this study.

\begin{table}[t]
\centering
\scriptsize
\resizebox{\linewidth}{!}{\input{generated/configs.tex}}
\caption{Configurations. Values are read from the release rows. The fx and fxh model id is the same model under its Vercel AI Gateway name. Differences inside the fixed-model condition are listed in Section~\ref{sec:fixed}.}
\label{tab:configs}
\end{table}

\subsection{Tasks}

A task is a Harbor task directory (instruction, container environment, oracle solution, verifier) with a sidecar that declares the poisonable resources, non-idempotent tools and alternatives.
Service-track tasks depend on mock services (tickets, inventory, deploy, documentation); native-track tasks are small repositories with bug fixes, features, builds, investigations, refactors and data tasks.
All tasks are written for the benchmark; none comes from a production project.
Every task passes per-task validation in continuous integration: the oracle scores 1 under the clean profile and, on the service track, under every fault profile; a null agent scores 0.
A task then enters the benchmark only after calibration with Claude Code under \code{clean} (5 attempts): clean pass$^1$ of at least 0.4 and a floor of eligible tool calls (8 on the service track, 4 on the native track), so that faults have calls to land on.
Because calibration uses one of the compared configurations, the task set can favor tasks that Claude Code solves cleanly; Section~\ref{sec:threats} discusses this.
The verifier checks final state deterministically and lives outside the agent's image.
Every agent-hidden task file carries a canary string.
Each instruction ends with a completion contract: the agent writes a status file stating \code{success}, \code{partial} or \code{failure}, which the overclaim metric reads.

\subsection{Profiles and attempts}

A profile sets the per-call base rate, the per-trial cap, the persistence mix, the type weights and the targeted rates (Table~\ref{tab:profiles}).
The headline profile is \code{moderate}.
Its type weights on the service track are \code{err} 0.20, \code{trunc} and \code{falsefail} 0.15 each, and \code{deny}, \code{empty}, \code{poison}, \code{drift} and \code{rate} 0.10 each, in every faulted profile; native-track weights are renormalized over the parity set.
The \code{clean} profile applies no faults and gives the baseline.
The profiles change several parameters together, so they are severity levels, not a manipulation of the base rate alone.
Each configuration runs every task under \code{clean} and \code{moderate} with $n = 3$ attempts, and attempt $a$ of every configuration uses the same schedule.

\begin{table}[t]
\centering
\footnotesize
\begin{tabular}{lrrll}
\toprule
Profile & Base rate & Cap per trial & Persistence & Targeted faults \\
\midrule
\code{clean} & 0 & 0 & none & none \\
\code{light} & 0.05 & 3 & transient & none \\
\code{moderate} & 0.15 & 6 & transient 70\%, \code{sticky(2)} 30\% & half of trials \\
\code{heavy} & 0.30 & 10 & transient 50\%, \code{sticky(2)} 35\%, outage 15\% & every trial \\
\bottomrule
\end{tabular}
\caption{Fault profiles, from \code{profiles/*.yaml}. An outage fails every call to one MCP tool for the next 3 eligible calls; it applies only on the service track and only to tools with a declared alternative. Targeted false failures apply only on the service track.}
\label{tab:profiles}
\end{table}
